\begin{figure}[!htbp]
  \centering
  \resizebox{\linewidth}{!}{%
  \begin{tikzpicture}[
    x=1cm, y=1cm,
    every node/.style={font=\scriptsize, align=center},
    panel/.style={
      rectangle, rounded corners=5pt,
      draw=black!25, fill=black!1, line width=0.7pt
    },
    title/.style={font=\bfseries\small, anchor=west, text=black},
    note/.style={font=\scriptsize, anchor=west, text=black!64},
    statebox/.style={
      rectangle, rounded corners=3pt,
      minimum width=0.72cm, minimum height=0.43cm,
      inner sep=2.1pt, line width=0.75pt
    },
    chainstep/.style={
      rectangle, rounded corners=3pt,
      text width=4.55cm, minimum height=0.48cm,
      inner sep=3.6pt, draw=black!38, fill=white, line width=0.65pt
    },
    arr/.style={-stealth, draw=black!58, line width=0.75pt}
  ]
    \definecolor{kBlue}{RGB}{58,104,173}
    \definecolor{kBlueFill}{RGB}{224,235,249}
    \definecolor{vTeal}{RGB}{66,139,133}
    \definecolor{vTealFill}{RGB}{224,242,239}
    \definecolor{riskRed}{RGB}{184,84,78}
    \definecolor{riskRedFill}{RGB}{249,228,225}
    \definecolor{anchorGold}{RGB}{178,130,26}
    \definecolor{anchorGoldFill}{RGB}{250,239,202}
    \definecolor{safeGreen}{RGB}{73,135,96}
    \definecolor{safeGreenFill}{RGB}{230,243,234}

    % panel backgrounds
    \node[panel, minimum width=8.65cm, minimum height=5.65cm, anchor=south west] at (0.00,0.00) {};
    \node[panel, minimum width=5.85cm, minimum height=5.65cm, anchor=south west] at (8.95,0.00) {};

    % panel (a)
    \node[title] at (0.30,5.30) {（a）位宽落点诊断};
    \node[note] at (0.30,4.95) {同平均位宽预算下，4bit 应优先落在 K 还是 V？};

    \node[font=\tiny\bfseries, text=black!70] at (1.85,4.42) {8 位基准};
    \draw[black!28, line width=0.55pt] (1.28,4.26) -- (2.42,4.26);
    \node[font=\tiny\bfseries, text=black!70] at (3.95,4.42) {同预算对照};
    \draw[black!28, line width=0.55pt] (2.75,4.26) -- (5.15,4.26);
    \node[font=\tiny\bfseries, text=black!70] at (6.25,4.42) {对称低比特锚点};
    \draw[black!28, line width=0.55pt] (5.60,4.26) -- (6.90,4.26);

    \node[font=\scriptsize\bfseries] at (1.85,3.94) {\texttt{K8V8}};
    \node[font=\scriptsize\bfseries] at (3.25,3.94) {\texttt{K8V4}};
    \node[font=\scriptsize\bfseries] at (4.65,3.94) {\texttt{K4V8}};
    \node[font=\scriptsize\bfseries] at (6.25,3.94) {\texttt{K4V4}};

    \node[anchor=east, font=\tiny\bfseries, align=right] at (1.05,3.22) {Qwen\\1.5B};
    \node[anchor=east, font=\tiny\bfseries, align=right] at (1.05,2.42) {Qwen\\7B};
    \node[anchor=east, font=\tiny\bfseries, align=right] at (1.05,1.62) {LLaMA\\8B};

    \foreach \y in {2.82,2.02,1.22} {
      \draw[black!10] (1.28,\y) -- (7.10,\y);
    }

    \node[statebox, fill=safeGreenFill, draw=safeGreen, text=safeGreen!65!black] at (1.85,3.22) {可用};
    \node[statebox, fill=safeGreenFill, draw=safeGreen, text=safeGreen!65!black] at (3.25,3.22) {可用};
    \node[statebox, fill=riskRedFill, draw=riskRed, text=riskRed!80!black] at (4.65,3.22) {崩塌};
    \node[statebox, fill=riskRedFill, draw=riskRed, text=riskRed!80!black] at (6.25,3.22) {崩塌};

    \node[statebox, fill=safeGreenFill, draw=safeGreen, text=safeGreen!65!black] at (1.85,2.42) {可用};
    \node[statebox, fill=safeGreenFill, draw=safeGreen, text=safeGreen!65!black] at (3.25,2.42) {可用};
    \node[statebox, fill=riskRedFill, draw=riskRed, text=riskRed!80!black] at (4.65,2.42) {崩塌};
    \node[statebox, fill=riskRedFill, draw=riskRed, text=riskRed!80!black] at (6.25,2.42) {崩塌};

    \node[statebox, fill=safeGreenFill, draw=safeGreen, text=safeGreen!65!black] at (1.85,1.62) {可用};
    \node[statebox, fill=safeGreenFill, draw=safeGreen, text=safeGreen!65!black] at (3.25,1.62) {可用};
    \node[statebox, fill=safeGreenFill, draw=safeGreen, text=safeGreen!65!black] at (4.65,1.62) {保留};
    \node[statebox, fill=safeGreenFill, draw=safeGreen, text=safeGreen!65!black] at (6.25,1.62) {保留};

    \node[font=\tiny, text=black!62, align=left, anchor=west] at (0.65,0.48)
      {代表性检索观察（Needle）：K8V8 与 K8V4 保持可用；\\Qwen 系列在 4bit K 下更易崩塌；LLaMA-3.1-8B 保留较高通过率，\\说明强度受架构调制。};

    % panel (b)
    \node[title] at (9.25,5.30) {（b）设计启示链};
    \node[note] at (9.25,4.95) {从诊断观察收束为低比特路径约束};

    \node[chainstep, fill=kBlueFill, draw=kBlue] (c1) at (11.88,4.28)
      {高精度路径保持可用};
    \node[chainstep, fill=vTealFill, draw=vTeal] (c2) at (11.88,3.34)
      {同预算落点暴露失效侧};
    \node[chainstep, fill=riskRedFill, draw=riskRed] (c3) at (11.88,2.40)
      {4bit K 可能越过注意力排序临界点};
    \node[chainstep, fill=anchorGoldFill, draw=anchorGold] (c4) at (11.88,1.42)
      {优先保护 K，再区分 K/V 角色};
    \node[chainstep, fill=black!4, draw=black!45] (c5) at (11.88,0.52)
      {角色感知非对称格式};

    \draw[arr] (c1.south) -- (c2.north);
    \draw[arr] (c2.south) -- (c3.north);
    \draw[arr] (c3.south) -- (c4.north);
    \draw[arr] (c4.south) -- (c5.north);
  \end{tikzpicture}%
  }
  \caption{K/V 位宽落点诊断的压缩版方法动机。子图（a）将 \texttt{K8V8} 基准、同平均位宽预算下的 \texttt{K8V4}/\texttt{K4V8} 对照，以及 \texttt{K4V4} 对称低比特锚点合并为 $3\times4$ 诊断矩阵，只保留代表性 Needle 观察：\texttt{K8V8} 与 \texttt{K8V4} 保持可用，Qwen 系列在 \texttt{K4V8}/\texttt{K4V4} 下出现 K 侧崩塌，而 LLaMA-3.1-8B 不呈现同强度坍塌。子图（b）将该观察收束为设计启示链：高精度路径保持可用，同预算落点暴露失效侧，4bit K 可能越过注意力排序临界点，因此低比特路径应优先保护 K，并进一步区分 K/V 角色。本图不承担完整实验比较；完整 Needle、RULER 与 LongBench 风格证据见第~\ref{subsec:ch4-kv-diagnosis}~节。}
  \label{fig:ch3-kv-asymmetry}
\end{figure}
